\section{Requisitos Funcionales}
Este capítulo describe de forma estructurada \textbf{las acciones que el sistema debe
ser capaz de realizar} para cumplir los objetivos establecidos en los capítulos
anteriores. 
\par
\vspace{4mm}
Cuando un requisito tiene reglas de negocio asociadas o
condiciones de error que merecen detalle, estas se listan bajo
la descripción principal.

\subsection{RF1 Gestión de Cuentas y Autenticación}
\begin{itemize}
    \item \textbf{RF1.1 Creación de cuentas de voluntario} \\
    El sistema debe permitir al usuario con rol Admin o voluntario con rol adecuado crear nuevas cuentas de voluntario en el sistema. 
    Los datos que se registran en el formulario son:
    \begin{description}
        \item[-] Nombre.
        \item[-] Apellidos.
        \item[-] Teléfono.
        \item[-] DNI. % chktex 13
        \item[-] Correo electrónico.
        \item[-] Rol.
        \item[-] Nivel de seguridad.
        \item[-] Area profesional del voluntario. 
        \item[-] Capacidades en el sistema.
    \end{description}
    \par
    Adicionalmente, el sistema registrará automáticamente la siguiente información de la cuenta en todo momento:
    \begin{description}
        \item[-] Estado de la cuenta.
        \item[-] Contraseña encriptada.
        \item[-] Fecha de creación de la cuenta al finalizar la creación.
        \item[-] Fecha del último login del voluntario.
        \item[-] Intentos fallidos al intentar iniciar sesión de forma incorrecta.
        \item[-] Fecha hasta donde se bloquea la cuenta tras 5 intentos fallidos al iniciar sesión. 
    \end{description}
    \par
    La creación de una nueva cuenta queda registrada en el log de auditoría con el nombre de la cuenta actualmente autenticada en el sistema. 
    \par
    El admin proporciona al voluntario una contraseña de un solo uso para entrar al sistema. Una vez inicie el voluntario, se le pedirá que introduzca 
    su nueva contraseña ed acceso para sustituir la temporal proporcionada. En este momento, el estado de la cuenta se actualiza de \texttt{PENDIENTE\_ACTIVACION} a 
    \texttt{ACTIVA} (RF1.8).
    
    \item \textbf{RF1.2 Visualización de los datos de la cuenta} \\
    El sistema debe permitir a cualquier voluntario autenticado consultar los
    datos de su propia cuenta.

    \item \textbf{RF1.3 Edición de los datos de la cuenta} \\
    El sistema debe permitir a cualquier voluntario autenticado modificar los
    datos personales de su propia cuenta.

    \item \textbf{RF1.4 Eliminación de cuentas} \\
    El sistema debe permitir al Admin eliminar definitivamente una cuenta de
    voluntario sólo si esta no tiene operaciones históricas asociadas.
    \begin{description}
        \item[-] Si la cuenta tiene operaciones registradas, la eliminación se rechaza
        y se ofrece la alternativa de desactivación (RF1.5).
        \item[-] La eliminación deja constancia en el log de auditoría con el
        identificador y nombre de la cuenta eliminada.
    \end{description}

    \item \textbf{RF1.5 Desactivación de cuentas} \\
    El sistema debe permitir al Admin desactivar la cuenta de un voluntario sin eliminarla del sistema.
    \begin{description}
        \item[-] La desactivación es reversible por el Admin en cualquier momento.
        \item[-] Una cuenta desactivada no permite inicio de sesión.
    \end{description}

    \item \textbf{RF1.6 Inicio de sesión} \\
    El sistema debe permitir a un voluntario iniciar sesión introduciendo su
    correo electrónico y contraseña.
    \begin{description}
        \item[-] Tras un inicio de sesión correcto, el sistema emite un token JWT con
        validez limitada que el cliente almacena en memoria del SPA.\@
        \item[-] Los intentos fallidos quedan registrados en el log de auditoría. Varios intentos fallidos 
        consecutivos provoca un bloqueo temporal el inicio de sesión.
    \end{description}

    \item \textbf{RF1.7 Cierre de sesión} \\
    El sistema debe permitir a un voluntario cerrar sesión de forma explícita invalidando el token JWT 
    del lado del cliente.
    \begin{description}
        \item[-] El sistema debe ofrecer cierre de sesión automático tras un periodo
        configurable de inactividad.
    \end{description}

    \item \textbf{RF1.8 Activación de la cuentas} \\
    El sistema debe permitir a un voluntario activar su cuenta actualizando la contraseña temporal de acceso por una nueva que introduzca. 
    El estado de cuenta debe ser \texttt{ACTIVA} tras esta acción.
\end{itemize}




\subsection{RF2 Control de Acceso y Roles}
\begin{itemize}
    \item \textbf{RF2.1 Roles del sistema} \\
    El sistema debe implementar los cinco roles definidos en la sección 4.3, cada uno
    con su matriz de permisos por defecto: \texttt{ADMIN}, \texttt{OPERATIVO},
    \texttt{PROFESIONAL}, \texttt{PUENTE} y \texttt{AGENTE\_TEMP}.

    \item \textbf{RF2.2 Asignación de funciones específicas} \\
    El sistema debe permitir al Admin asignar a un voluntario, además de su
    rol, un conjunto de funciones específicas que limiten o amplíen los
    permisos por defecto del rol.
    \begin{description}
        \item[-] Las funciones específicas no pueden otorgar a un voluntario
        capacidades que su rol no le permita por diseño (un \texttt{OPERATIVO} no
        puede recibir la capacidad \texttt{IMPORTAR\_EXTRACTOS}, exclusiva del Admin).
    \end{description}

    \item \textbf{RF2.3 Niveles de visibilidad de datos sensibles} \\
    El sistema debe asignar a cada voluntario uno de los niveles de
    seguridad descritos en la sección 4.3: \texttt{BASICO} o \texttt{PRIVILEGIADO}.
    \begin{description}
        \item[-] El nivel \texttt{BASICO} ve únicamente los datos no sensibles de la ficha del amigo y
        los servicios prestados a él más no el motivo o notas que justifiquen esos servicios.
        \item[-] El nivel \texttt{PRIVILEGIADO} ve todos los datos. El Admin tiene siempre nivel y puede otorgarlo a otros voluntarios.
    \end{description}

    \item \textbf{RF2.4 Habilitación temporal de acceso} \\
    El sistema debe permitir al Admin habilitar a un voluntario profesional
    acceso de lectura a los datos sensibles de un amigo concreto durante una
    ventana de tiempo limitada.
    \begin{description}
        \item[-] La habilitación se asocia al par voluntario-amigo y a una fecha de
        caducidad. Alcanzada la fecha de caducidad, el acceso se revoca automáticamente sin intervención manual.
        \item[-] Toda actividad realizada bajo una habilitación temporal queda
        registrada en el log de auditoría con marcado especial.
    \end{description}
    
    \item \textbf{RF2.5 Revocación de habilitaciones temporales} \\
    El sistema debe permitir al Admin revocar una habilitación temporal antes
    de su caducidad.
    \begin{description}
        \item[-] La revocación es inmediata: la siguiente petición del voluntario
        afectado falla con código de autorización denegada.
        \item[-] La revocación queda registrada en el log de auditoría.
    \end{description}
    
    % FIXME: No se interpretar este requisito. ¿Qué es exactamente la "verificación de autorización en backend"?
    \item \textbf{RF2.6 Verificación de autorización en backend} \\
    El sistema debe verificar las reglas de autorización en el lado del
    servidor para toda operación que acceda a datos sensibles o que
    modifique el estado del sistema.
\end{itemize}




\subsection{RF3 Gestión de Amigo}
\begin{itemize}
    \item \textbf{RF3.1 Creación de ficha mínima} \\
    El sistema debe permitir crear una ficha de amigo introduciendo
    únicamente nombre.
    \begin{description}
        \item[-] El cambio de nivel queda registrado en el log de auditoría con 
        identificador del voluntario que lo realiza y timestamp.
        \item[-] La creación de la ficha queda registrada en el log de auditoría.
        \item[-] Se registra la fecha de creación en su ficha. 
    \end{description}
    \begin{itemize}
        \item \textbf{RF3.1.1 Registro de entrega objetos} \\
        En cualquier nivel de confianza (RF3.2) está disponible una caja de texto donde se registran objetos o pertenencias 
        del amigo al momento de entrar en el centro. Si son más de uno deben separarse por comas (, ) para que el sistema los procese % chktex 37
        correctamente y se muestren en la ficha del amigo en forma de listado de checkboxes. Caso paradigmático: amigo entrega pasaporte/DNI para que se guarde en la caja 
        fuerte del centro. En su lugar, se les otorga una copia.
        \begin{description}
            \item[-] El texto escrito aparecerá en la ficha correspondiente del amigo (RF3.4). Puede dejarse en blanco.
        \end{description}
        \item \textbf{RF3.1.2 Registro de devolución de objetos} \\
        El sistema registra en el log de auditoría la devolución de un objeto a un amigo con el nombre del usuario actualmente autenticado que realizó la devolución. Se actualiza consecuentemente su ficha de amigo 
        para reflejar que el objeto ya no está en posesión del centro. Caso paradigmático: se devuelve el DNI original al amigo para realizar un trámite legal.
    \end{itemize}

    \item \textbf{RF3.2 Escalado del nivel de confianza}
    \begin{itemize}
        \item \textbf{RF3.2.1 Nivel 1 $-$ Anónimo} \\
        Estado inicial. Campos disponibles: nombre y, opcionalmente, una descripción corta.
        \item \textbf{RF3.2.2 Nivel 2 $-$ Frecuente} \\
        Tras varias visitas o cuando la coordinación lo decida. Habilita los
        campos como pueden ser: apellidos, fecha de nacimiento, género, tallas (ropa y
        calzado), idioma principal, teléfono.
        \item \textbf{RF3.2.3 Nivel 3 $-$ Entrevistado} \\
        Tras una entrevista de profundidad. Habilita los campos: DNI o
        documento equivalente (NIE, pasaporte), nacionalidad, dirección, datos
        sensibles (situación administrativa, médica, psicológica, jurídica),
        notas privadas.
    \end{itemize}
    
    \item \textbf{RF3.3 Edición de ficha} \\
    El sistema debe permitir editar los campos disponibles en el nivel actual de la ficha.
    \begin{description}
        \item[-] Cada modificación queda registrada en el log de auditoría con valor 
        anterior, valor nuevo, voluntario y marca de tiempo.
        \item[-] Todos los campos son editables por el Admin o voluntariado con permisos adecuados.
    \end{description}
    
    \item \textbf{RF3.4 Visualización de ficha} \\
    El sistema debe permitir visualizar la ficha de un amigo aplicando las
    reglas de visibilidad por nivel de seguridad descritas en RF2.3.
    
    \item \textbf{RF3.5 Búsqueda y listado de amigos} \\
    El sistema debe permitir buscar y listar amigos aplicando filtros por
    nombre, alias, nivel de confianza y fecha de última visita mostrando información acorde a
    el nivel de seguridad del voluntario que consulta.
    
    
    \item \textbf{RF3.6 Histórico individual del amigo}\hypertarget{RF3.6}{} \\
    El sistema debe mantener y mostrar para cada amigo el histórico
    completo de su relación con la asociación:
    \begin{description}
        \item[-] Servicios solicitados, agendados y prestados (con fecha, voluntario, observaciones).
        \item[-] Entregas materiales recibidas (prenda, talla, fecha, voluntario).
        \item[-] Operaciones de lavandería (entrada de prenda propia, devolución).
        \item[-] Entradas y salidas alimentarias asociadas (cuando aplique).
        \item[-] Gastos vinculados al amigo (RF7.6).
        \item[-] Comunicaciones externas registradas (RF3.8).
        \item[-] Cambios de nivel de confianza con fechas y motivos.
    \end{description}

    \item \textbf{RF3.7 Registro de tallas} \\
    El sistema debe almacenar las tallas de ropa y calzado del amigo.
    
    \item \textbf{RF3.8 Trazabilidad de comunicaciones externas} \\
    El sistema debe permitir registrar manualmente comunicaciones externas
    realizadas en nombre del amigo o relacionadas con su expediente, con el
    fin de mantener un historial consolidado.
    \begin{description}
        \item[-] Tipos de comunicación: correo electrónico, llamada telefónica, gestión presencial, documento enviado o recibido.
        \item[-] Cada comunicación se registra con: tipo, contraparte (entidad o persona externa), fecha, asunto y notas opcionales.
    \end{description}

    \item \textbf{RF3.9 Notas privadas} \\
    El sistema debe permitir asociar notas privadas a la ficha de un amigo. Son
    visibles únicamente para voluntarios con nivel de seguridad \texttt{PRIVILEGIADO} o
    con habilitación temporal vigente al área correspondiente más cualquier voluntario puede añadir notas. 
    Caso paradigmático: un amigo le cuenta su historia por partes a diferentes voluntarios (con los que más simpatiza será más honesto). Registran esas notas puntuales. 
    Al cabo de varias notas se puede armar una linea temporal completa de la vida del amigo.
    \begin{description}
        \item[-] Las notas son inmutables una vez guardadas: no se editan, sino que se
        añaden notas adicionales que matizan o corrigen las previas.
        \item[-] Cada nota lleva identificador del voluntario autor y un timestamp de creación.
    \end{description}
\end{itemize}




\subsection{RF4 Servicios y Agendamiento}
\begin{itemize}
    \item \textbf{RF4.1 Gestión del Catálogo de servicios} \\
    El sistema debe mantener un catálogo de tipos de servicio que el Admin
    puede consultar y eliminar o desactivar algún según sea el caso.
    \begin{description}
        \item[-] La eliminación de un tipo de servicio queda condicionada a la ausencia
        de servicios prestados de ese tipo en el histórico (en su lugar, se
        desactiva).
    \end{description}

    \item \textbf{RF4.2 Agendamiento de un servicio} \\
    El sistema debe permitir agendar la prestación de un servicio a un amigo en una fecha y franja horaria concretas.
    \begin{description}
        \item[-] El sistema verifica al agendar que no se excede la capacidad del turno (RF4.3).
    \end{description}

    \item \textbf{RF4.3 Control de capacidad por turno} \\
    El sistema debe impedir el agendamiento de un servicio que excedería la
    capacidad máxima del turno correspondiente.
    
    \item \textbf{RF4.4 Regla de prerrequisito de agendamiento} \\
    El sistema debe impedir el registro de prestación de un servicio si no existe un agendamiento previo asociado.
    \begin{description}
        \item[-] La ducha es el caso paradigmático: sin agendamiento previo no se presta el servicio aunque haya espacio disponible.
    \end{description}
    
    \item \textbf{RF4.5 Registro de prestación de servicio efectiva} \\
    El sistema debe permitir al voluntario registrar la prestación efectiva
    del servicio en el momento o tras producirse.
    \begin{description}
        \item[-] Resultado posibles: prestado, no prestado por ausencia del amigo, no
        prestado por ausencia del voluntario, cancelado por la coordinación.
        \item[-] En caso de prestación, se registran observaciones opcionales del motivo de prestación del servicio 
        sólo visible para voluntarios autorizados.
        \item[-] La prestación queda asociada al amigo y aparece en su histórico (RF3.6).
    \end{description}
    
    \item \textbf{RF4.6 Visualización de la agenda} \\
    El sistema debe permitir visualizar la agenda del centro en formato calendario diario, semanal y mensual.
    \begin{description}
        \item[-] Cada agendamiento se muestra con tipo de servicio, amigo asociado y voluntario responsable.
        \item[-] La vista diaria es accesible directamente desde la pantalla principal.
        \item[-] Filtros disponibles: por tipo de servicio, por voluntario, por amigo.
    \end{description}
    
    \item \textbf{RF4.7 Reagendamiento y cancelación} \\
    El sistema debe permitir modificar la fecha, franja o voluntario asignado de un agendamiento, así como cancelarlo.
    \begin{description}
        \item[-] Cualquier modificación queda registrada en el log de auditoría.
        \item[-] La cancelación no elimina el registro: lo marca como cancelado con motivo opcional y voluntario que cancela.
    \end{description}
\end{itemize}




\subsection{RF5 Inventario, Entregas y Lavandería}
\begin{itemize}
    \item \textbf{RF5.1 Cuaderno de Entradas y Salidas} \\
    El sistema debe mantener un inventario del stock disponible en el
    armario general, organizado por tipo de prenda y talla.
    \begin{description}
        \item[-] El sistema muestra el inventario en una vista de cuadrícula que
        permite filtrar por tipo y talla. 
    \end{description}

    \item \textbf{RF5.2 Registro de entrada de stock} \\
    El sistema debe permitir registrar la entrada de prendas donadas al
    armario general.
    \begin{description}
        \item[-] La entrada incrementa automáticamente el stock correspondiente.
        \item[-] Es posible registrar varias prendas en una misma operación (entrada por lotes). 
    \end{description}

    \item \textbf{RF5.3 Registro de entrega de prenda} \\
    El sistema debe permitir registrar la entrega de una prenda concreta a un amigo.
    
    \begin{itemize}
        \item \textbf{RF5.3.1 Decremento automático de stock} \\
        La entrega decrementa automáticamente la cantidad disponible de la
        combinación tipo-talla entregada en una unidad.
        \item \textbf{RF5.3.2 Reglas de frecuencia configurables} \\
        El sistema debe aplicar reglas de frecuencia configurables que limiten
        la entrega repetida del mismo tipo de prenda al mismo amigo en un
        periodo dado.
        \begin{description}
            \item[-] Caso paradigmático: una unidad de calzado por amigo cada 3 meses.
            \item[-] Si la regla se incumple, el sistema avisa al voluntario y exige justificación o aprobación de un Admin.
            \item[-] Las excepciones quedan registradas en el log de auditoría. 
        \end{description}
        \item \textbf{RF5.3.3 Integridad ante stock insuficiente} \\
        El sistema debe impedir la entrega de una prenda cuya cantidad disponible es cero. 
        Ni siquiera debe mostrarse como opción seleccionable en el proceso de entrega.
    \end{itemize}

    \item \textbf{RF5.4 Alertas de reposición} \\
    El sistema debe permitir definir, por combinación tipo-talla, un umbral
    mínimo de stock por debajo del cual se genera una alerta visual en el
    panel del armario.

    \item \textbf{RF5.5 Visualización del histórico de lavandería} \\
    El sistema debe ofrecer una vista del histórico de lavandería tanto a
    nivel global del centro (todas las prendas en proceso) como a nivel
    individual del amigo.
    \begin{description}
        \item[-] El histórico individual aparece en la ficha del amigo (RF3.6).
        \item[-] El histórico global es accesible desde el módulo de inventario.
    \end{description}
\end{itemize}




\subsection{RF6 Trazabilidad Alimentaria (Cumplimiento Ley 1/2025)}
\begin{itemize}
    \item \textbf{RF6.1 Registro de entrada de alimento} \\
    El sistema debe permitir registrar la entrada de alimentos al centro,
    con todos los datos necesarios para satisfacer las exigencias de
    trazabilidad de la Ley 1/2025.
    \begin{itemize}
        \item \textbf{RF6.1.1 Origen del alimento} \\
        Cada entrada debe categorizarse según su origen:
        \begin{description}
            \item[-] Compra propia: realizada por la asociación. Requiere ticket de compra asociado.
            % \item[-] Donación de entidad con convenio: requiere convenio activo registrado (RF6.x) y, cuando proceda, albarán de entrega. % No se incluye en esta fase 
            \item[-] Donación de entidad sin convenio: para donaciones puntuales de entidades sin acuerdo formal.
            \item[-] Donación particular: Persona que aporta espontáneamente. En este caso, el ticket suele no estar disponible y debe quedar documentado el motivo.
        \end{description}
        \par 
        Todo ticket o albarán recibido debe almacenarse físicamente y estar disponible para auditoría.
        \item \textbf{RF6.1.2 Datos obligatorios de la entrada} \\
        Independientemente del origen, toda entrada registra: 
        \begin{description}
            \item[-] Tipo de producto.
            \item[-] Categoría general (lácteos, panadería, fruta\ldots)
            \item[-] Cantidad y unidad (kilos, litros, unidades).
            \item[-] Fecha de entrada.
            \item[-] Voluntario que registra.
            \item[-] Lugar de almacenamiento (despensa, congelador\ldots)
        \end{description}
        \item \textbf{RF6.1.3 Datos específicos de seguridad alimentaria} \\
        Para todo producto perecedero o susceptible de control sanitario:
        \begin{description}
            \item[-] Lote del fabricante, cuando esté disponible.
            \item[-] Fecha de caducidad o consumo preferente.
            \item[-] Para productos refrigerados: temperatura de recepción.
            \item[-] Check\footnote{Se refiere a una verificación física, y se marca con un \textit{check} que el estado del producto es aceptable.} de verificación de higiene y estado del producto al recibirse.
            \item[-] Lugar y método de almacenamiento.
        \end{description}
    \end{itemize}

    \item \textbf{RF6.2 Registro de salida de alimento} \\
    El sistema debe permitir registrar la salida de alimento del stock
    hacia su consumo o destino final.
    \begin{description}
        \item[-] Cada salida incluye: producto, cantidad, fecha, turno (mañana o tarde), tipo de salida y voluntario que registra.
        \item[-] Tipos de salida: desayuno, merienda, cena, café, comida individual a un amigo, retirada por caducidad, retirada por deterioro.
        \item[-] La salida decrementa automáticamente el stock disponible.
        \item[-] Se registra el voluntario que realizo la salida.
    \end{description}

    \item \textbf{RF6.3 Clasificación por prioridad legal} \\
    Toda entrada y, cuando aplique, toda salida del sistema deben
    clasificarse según las cuatro categorías de prioridad de aprovechamiento
    establecidas por la Ley 1/2025:
    \begin{description}
        \item[-] Prioridad 1: Consumo humano. Categoría por defecto.
        \item[-] Prioridad 2: Alimentación animal.
        \item[-] Prioridad 3: Subproducto (compostaje, biocombustible).
        \item[-] Prioridad 4: Residuo.
        \item[-] Toda salida con prioridad distinta de 1 requiere justificación
        obligatoria: motivo (caducidad, deterioro, rechazo, otros) y
        observaciones.
        \item[-] El sistema mantiene un contador agregado por prioridad para los informes legales.
    \end{description}

    \item \textbf{RF6.4 Control de stock alimentario} \\
    El sistema debe mantener una vista de stock alimentario actual
    diferenciada del armario general (RF5.1), con visibilidad de:
    \begin{description}
        \item[-] Cantidad disponible por tipo de producto.
        \item[-] Lotes activos con sus fechas de caducidad.
        \item[-] Productos próximos a caducar (umbral configurable por defecto a 7 días).
        \item[-] Productos caducados pendientes de retirada.
    \end{description}

    % \item \textbf{RF6.5 Asociación de salida a amigo} \\
    % El sistema debe permitir, opcionalmente, asociar una salida concreta a un amigo identificado.
    % \begin{description}
    %     \item[-] Esta asociación puede incrementar la carga de trabajo de los voluntarios en el día a día. Sin embargo, se 
    %     considera que refuerza la trazabilidad y permite generar
    %     informes de consumo individualizado (por ejemplo, \textit{``café tomado por
    %     Marcos durante el último semestre''}).
    %     \item[-] La asociación, cuando se usa, aparece en el histórico del amigo (RF3.6). 
    % \end{description}

    \item \textbf{RF6.5 Generación del cuaderno de trazabilidad} \\
    El sistema debe generar el cuaderno completo de entradas y
    salidas alimentarias en un periodo definido, en formato apto para
    inspección sanitaria.
    \begin{description}
        \item[-] Formatos aptos: PDF (impresión) y CSV (análisis).
        \item[-] Contenido: todas las entradas con sus datos completos (RF6.1), todas
        las salidas (RF6.2), todos los productos retirados con justificación
        (RF6.3) y saldo de stock al cierre del periodo.
        \item[-] El cuaderno generado para un periodo cerrado debe ser inmutable: no
        puede regenerarse con datos modificados a posteriori sin que quede
        constancia de la modificación.
    \end{description}

    \item \textbf{RF6.6 Inmutabilidad de registros cerrados} \\
    El sistema debe impedir la modificación o eliminación de registros de
    entrada o salida con antigüedad superior a un umbral configurable (por
    defecto, 30 días).
    \begin{description}
        \item[-] Esta regla protege la integridad documental del cuaderno frente a inspecciones.
        \item[-] Si una corrección es necesaria tras el cierre, se realiza mediante asientos
        compensatorios (entrada o salida correctiva con motivo \texttt{Rectificación}), no mediante edición del registro original.
    \end{description}
\end{itemize}

\subsection*{RF7 Gestión Financiera y Trazabilidad de Gasto}
\begin{itemize}
    \item \textbf{RF7.1 Catálogo de cuentas bancarias y caja} \\
    El sistema debe permitir registrar y mantener las cuentas a través de
    las cuales la asociación gestiona sus movimientos.
    \begin{description}
        \item[-] Para cada cuenta bancaria: nombre, entidad, número de cuenta 
        parcialmente enmascarado, tipo (corriente, caja chica, otra), estado (activa, inactiva).
        \item[-] La caja chica se modela como una cuenta de tipo efectivo.
        \item[-] Una cuenta no puede eliminarse si tiene asientos asociados; en su lugar, se desactiva. 
    \end{description}

    \item \textbf{RF7.2 Catálogo de categorías de gasto e ingreso} \\
    El sistema debe mantener un catálogo de categorías funcionales que el
    Admin puede ampliar y modificar.
    \begin{description}
        \item[-] El catálogo inicial reproduce las categorías que la asociación utiliza actualmente. Ejemplos: Suministros $-$ Agua/Luz, Gastos
        bancarios, Donaciones recibidas, Tasas de DNI, Medicamentos, productos de higiene, etc.
        \item[-] La eliminación de una categoría queda condicionada a la ausencia de asientos asociados; en su lugar, se desactiva.
    \end{description}

    \item \textbf{RF7.3 Catálogo de cuentas contables} \\
    El sistema debe permitir mantener, opcionalmente, un catálogo de
    códigos de cuenta contable conformes al criterio que la gestoría
    utiliza con la asociación.
\begin{description}
    \item[-] Cada cuenta contable incluye: código (por ejemplo ``6280'') y nombre descriptivo (por ejemplo ``Suministros $-$ Agua'').
    \item[-] El uso de cuenta contable en cada asiento es opcional: si la contabilidad llevada por la directiva no lo requiere, este catálogo puede permanecer vacío.
\end{description}

    \item \textbf{RF7.4 Registro manual de asientos contables} \\
    El sistema debe ofrecer un formulario para registrar manualmente cada
    asiento contable, reproduciendo digitalmente el flujo que la
    contabilidad de la asociación lleva hoy en Excel.
    \begin{itemize}
        \item \textbf{RF7.4.1 Campos del asiento} \\
        \begin{description}
            \item[-] \texttt{Fecha:} fecha del movimiento. Obligatorio.
            \item[-] \texttt{Concepto}: descripción libre del movimiento. Obligatorio.
            \item[-] \texttt{Cuenta bancaria}: cuenta o caja a la que se carga o abona el movimiento (RF7.1). Obligatorio.
            \item[-] \texttt{Cuenta contable}: código contable asociado (RF7.3). Opcional.
            \item[-] \texttt{Categoría}: clasificación funcional del movimiento (RF7.2). Obligatorio.
            \item[-] \texttt{Importe}: cantidad en euros, con signo (negativo para gasto, positivo para ingreso). Obligatorio.
            \item[-] \texttt{Tipo de gasto}: Casa, Amigo o Mixto (RF7.4.4). Obligatorio para gastos; no aplica a ingresos. 
            \item[-] \texttt{Amigo vinculado}: ficha del amigo asociado al gasto, sólo habilitado si el tipo es Amigo o Mixto. Opcional pero recomendado.
            \item[-] \texttt{Servicio vinculado}: servicio del catálogo (RF4.1) al que se asocia el gasto. Opcional.
            \item[-] \texttt{Referencia al archivo físico}: texto libre donde el usuario puede anotar la ubicación del justificante físico (``Carpeta abril 2026, ticket 7''). Opcional.
        \end{description}
        \item \textbf{RF7.4.2 Validación de los datos} \\
        El sistema debe validar al envío del formulario:
        \begin{description}
            \item[-] Los campos obligatorios están cumplimentados.
            \item[-] El importe es un número con dos decimales máximo.
            \item[-] Las cuentas y categorías referenciadas existen y están activas.
        \end{description}
        \item \textbf{RF7.4.3 Persistencia y auditoría} \\
        Cada asiento creado queda registrado con:
        \begin{description}
            \item[-] Identificador interno único e inmutable.
            \item[-] Marca de tiempo de creación.
            \item[-] Voluntario que lo registró.
            \item[-] Las modificaciones posteriores quedan registradas en el log de
            auditoría con valor anterior, valor nuevo, voluntario y marca de
            tiempo (RNF de auditoría).
        \end{description}
        \item \textbf{RF7.4.4 Categorización del gasto por destinatario} \\
        Cada gasto se clasifica obligatoriamente en una de tres categorías de destinatario:
        \begin{description}
            \item[-] \texttt{Casa}: gasto puro de la asociación (suministros generales, servicios, insumos comunes).
            \item[-] \texttt{Amigo}: gasto trazable a un amigo concreto (tasas de DNI, medicamentos prescritos, pasajes individuales).
            \item[-] \texttt{Mixto}: compra que beneficia simultáneamente a la asociación y sus miembros.
        \end{description}
        \par
        Esta categorización es independiente de la categoría funcional (RF7.2)
        y permite generar de forma natural los informes de gasto por persona
        solicitados por la directiva (campos de amigo asociado RF7.4.1).
    \end{itemize}

    \item \textbf{RF7.5 Búsqueda en historial asientos previos y optimización de rellenado} \\
    El sistema debe permitir, dentro del formulario de creación de
    asientos, buscar asientos previos similares y utilizar uno de ellos
    como plantilla para el nuevo asiento para optimizar el proceso de registro.
    \begin{description}
        \item[-] El usuario introduce un término de búsqueda y el sistema devuelve los asientos previos cuyo concepto contiene el
        término, ordenados por fecha descendente.
        \item[-] El usuario selecciona un asiento de los resultados y pulsa ``Usar
        como plantilla'': el formulario se rellena con todos los campos relevantes del asiento seleccionado.
        \item[-] El asiento original utilizado como plantilla no se modifica: el resultado es siempre un asiento nuevo e independiente.
    \end{description}


    \item \textbf{RF7.6 Autocompletado y prevención de duplicados en campos catalogados} \\
    Los campos del formulario que apuntan a entidades catalogadas (cuenta
    bancaria, cuenta contable, categoría, amigo, servicio) deben ofrecer
    autocompletado y verificación de duplicados.
    \begin{itemize}
        \item \textbf{RF7.6.1 Creación de nuevas entradas desde el formulario} \\
        Para los campos donde está permitido (cuenta bancaria, cuenta
        contable, categoría), si el valor introducido no existe en el
        catálogo, el sistema permite crearlo desde el propio formulario sin
        abandonar la operación de creación de asiento.
        \par
        Para los campos cerrados (amigo, servicio), no se permite la
        creación desde aquí: la entrada debe existir previamente en su
        pantalla de gestión correspondiente.
        \item \textbf{RF7.6.2 Detección de posibles duplicados} \\
        Antes de persistir un valor nuevo en un catálogo, el sistema debe
        verificar si el valor introducido es similar a alguno existente y,
        en caso afirmativo, solicitar confirmación al usuario.
    \end{itemize}
    

    \item \textbf{RF7.7 Listado y consulta de asientos contables} \\
    El sistema debe ofrecer una vista de listado de los asientos
    contables registrados, con capacidades de filtrado por cualquier campo del asiento 
    y ordenación por cualquier columna del listado.
    \begin{description}
        \item[-] Paginación de resultados obligatoria.
    \end{description}

    \item \textbf{RF7.8 Edición y anulación de asientos contables} \\
    El sistema debe permitir editar y anular asientos contables dentro de
    un periodo configurable (30 días por defecto desde la creación); pasado ese periodo, las correcciones se
    realizan mediante un asiento compensatorio.


    \item \textbf{RF7.9 Gestión de caja chica} \\
    El sistema debe mantener un saldo informativo de la caja chica
    basado en los asientos asociados a la cuenta de tipo efectivo
    (RF7.1).
    \begin{description}
        \item[-] El saldo se actualiza automáticamente con cada asiento.
        \item[-] El sistema ofrece una pantalla específica de ``Caja chica'' donde el 
                Admin puede consultar el saldo actual, el listado de movimientos 
                recientes y registrar el cuadre periódico (comparación entre saldo 
                del sistema y efectivo realmente custodiado).
        \item[-] Las diferencias detectadas en cuadres se registran como asientos de ajuste con motivo explícito.
    \end{description}


    \item \textbf{RF7.10 Consultas e informes de trazabilidad} \\
    El sistema debe permitir consultar los asientos contables agregados
    desde múltiples ángulos, generando los indicadores que la directiva
    necesita para la gestión interna y la justificación a entidades
    financiadoras (ligado a RF 8.2). Posibles consultas son:
    \begin{description}
        % TODO: generar una consulta JPQL concreta para cada caso + una completamente libre a partir de filtros por campos del asiento?
        \item[-] \textbf{Por amigo}: total invertido en cada amigo en un periodo dado,
                con desglose por categoría. Responde a la pregunta operativa
                ``\textit{¿cuánto hemos gastado en Marcos este semestre?}''.
        \item[-] \textbf{Por categoría funcional}: total por categoría en un periodo
                dado (por ejemplo, ``gasto total en Suministros $-$ Agua en 2026'').
        \item[-] \textbf{Por servicio}: total vinculado a cada área de servicio (gasto
                vinculado a podología, jurídico, etc.).
        \item[-] \textbf{Por voluntario}: total de asientos registrados por cada
                voluntario, como mecanismo de control interno.
        \item[-] \textbf{Por periodo}: comparativa mensual, semestral y anual con
                evolución frente al periodo equivalente anterior.
        \item[-] \textbf{Saldo por cuenta bancaria}: ingresos menos gastos, total y desglose por mes.
    \end{description}


\end{itemize}


\subsection{RF8 Reportes, Informes y Visualización}
\begin{itemize}
    \item \textbf{RF8.1 Panel de indicadores} \\
    El sistema debe ofrecer una pantalla principal personalizada según el
    rol del voluntario, que muestra los indicadores más relevantes para su
    trabajo.
    \begin{description}
        \item[-] Para Admin y directiva: total de amigos activos, servicios prestados
        en el último mes, gasto del mes en curso, alertas pendientes de
        reposición de stock.
        \item[-] Para Voluntario operativo: agenda del día, stock con alertas, tareas asignadas.
        \item[-] Los indicadores se calculan con datos en tiempo real, no con capturas periódicas.
    \end{description}


    \item \textbf{RF8.2 Generación de informes numéricos} \\
    El sistema debe generar informes con cifras numéricas exactas (no porcentajes ni aproximaciones).
    \begin{description}
        \item[-] Tipo de informe configurable: actividad general, gasto por categoría,
        servicios prestados, atención individualizada, trazabilidad
        alimentaria. 
        \item[-] Periodo configurable: mes, trimestre, semestre, año, rango libre.
    \end{description}


    \item \textbf{RF8.3 Informes para entidades financiadoras} \\
    El sistema debe ofrecer plantillas de informe específicas para las
    entidades financiadoras que solicitan justificación de subvención.


    \item \textbf{RF8.4 Informe agregado a nivel de ``universo''} \\
    El sistema debe generar el informe agregado de actividad de la
    asociación en un periodo dado, contemplando la visión de ``universo''
    solicitada explícitamente por la coordinación: visión global de toda la
    actividad sin tener que reconstruir información a partir de fuentes
    dispersas.
    \begin{description}
        \item[-] Cifras totales de personas atendidas, servicios prestados, alimentos
        gestionados, ropa entregada, gasto ejecutado. 
        \item[-] Desglose por tipo de actividad y por mes.
        \item[-] Indicadores de evolución frente al periodo anterior. 
    \end{description}


    \item \textbf{RF8.5 Exportación de informes} \\
    Todo informe generado debe poder exportarse en los siguientes formatos:
    \begin{description}
        \item[-] PDF:\@ para presentación a entidades externas e impresión.
        \item[-] CSV:\@ para análisis adicional en hoja de cálculo.
        \item[-] JSON:\@ para integraciones futuras.
        \item[-] Las exportaciones de datos sensibles (datos médicos, jurídicos,
        notas privadas) sólo están disponibles para voluntarios con nivel de
        seguridad \texttt{PRIVILEGIADO} y quedan registradas en el log de auditoría.
    \end{description}
\end{itemize}

\subsection{RF9 Puente Analógico-Digital y Carga Inicial de Datos}
\begin{itemize}
    \item \textbf{RF9.1 Impresión de planillas} \\
    El sistema debe generar planillas imprimibles cuyos campos coinciden
    exactamente con los de las entidades digitales correspondientes. Cada plantilla tendrá un 
    identificador único que permita su asociación inequívoca a la entidad digital en el sistema.


    \item \textbf{RF9.2 Control de Stock de Ropa} \\
    El sistema debe ofrecer una pantalla específica de \texttt{Carga analógica} en
    la que un voluntario introduce los datos contenidos en
    planillas cumplimentadas a mano.
    \begin{description}
        \item[-] La carga debe incluir el identificador de la planilla impresa, lo que
        permite trazar qué papel originó qué registro digital.
        \item[-] Los registros cargados de esta manera quedan marcados como \texttt{Origen analógico}
        para diferenciarlos de los introducidos directamente en sesión, lo
        que permite identificarlos en informes y auditorías.
    \end{description}


    \item \textbf{RF9.3 Importación de datos históricos} \\
    El sistema debe permitir importar los datos históricos del centro desde las hojas Excel actuales que mantienen la coordinación.
    \begin{description}
        \item[-] Cada sección de la aplicación contará con una lógica de mapeo específica a los campos de la entidad digital destino. 
    \end{description}


    \item \textbf{RF9.4 Validación de la importación} \\
    El sistema debe ofrecer una validación previa a toda importación
    masiva, mostrando:
    \begin{description}
        \item[-] Número de registros (filas del excel) que se importarán correctamente. 
        \item[-] Registros con errores de formato (con detalle del error por fila). 
        \item[-] Resumen estadístico (totales, rangos de fechas). 
        \item[-] El usuario decide tras la validación si confirma la importación,
        cancela, o aplica importación parcial sólo de las filas válidas. 
    \end{description}


    \item \textbf{RF9.5 Reversión de importaciones} \\
    El sistema debe permitir, durante un periodo limitado tras la ejecución
    (por defecto 24 horas), revertir una importación completa en caso de
    detectarse problemas no advertidos durante la validación.
    \begin{description}
        \item[-] Si se aplica una reversión, quedará registrada en el log de auditoría. 
        \item[-] Pasado el periodo, la corrección de datos importados erróneamente
        debe hacerse caso por caso.
    \end{description}
\end{itemize}
